\section{A uniform simplex estimate and a boundary crossover}
\label{sec:frontier-boundary}

The compact-interior assumption in Theorem~\ref{thm:simplex-interior}
excludes bins approaching a lower-dimensional face. We now remove that
assumption from the Bessel estimate and study the transition in which some
positive coordinates become small. The transition below needs at least two
coordinates of order $N$. With only one such coordinate the persistence
scale is different, and we treat that case at the end of this section and
in Section~\ref{sec:vertex-completion}.

Put
\[
\Phi(x)=\sum_{j\ge0}\frac{x^j}{j!(j+1)!}
=\frac{I_1(2\sqrt{x})}{\sqrt{x}},\qquad \Phi(0)=1.
\]
As before, $S_n(u)$ is the probability of the positive count vector $n$
divided by the probability of one vertex, and $u=r/p$. None of the
statements depend on $d$ once the support is fixed.

\begin{lemma}[a uniform Laguerre--Bessel inequality]
\label{lem:uniform-laguerre-bessel}
For every integer $a\ge1$ and real $w>0$,
\begin{equation}\label{eq:uniform-laguerre-bessel}
0\le1-\frac{B_a(w)}{w\Phi(aw)}\le\frac w2,
\qquad
B_a(w)=\sum_{m=1}^{a}\binom{a-1}{m-1}\frac{w^m}{m!}.
\end{equation}
\end{lemma}

\begin{proof}
The coefficient of $(aw)^j/[j!(j+1)!]$ in $B_a(w)/w$ is
$\prod_{h=1}^j(1-h/a)_+$. Its deficit from one lies between zero and
$j(j+1)/(2a)$, including for $j\ge a$. Consequently
\[
0\le\Phi(aw)-\frac{B_a(w)}w
\le\frac1{2a}\sum_{j\ge0}\frac{j(j+1)(aw)^j}{j!(j+1)!}
=\frac w2\Phi(aw).
\]
\end{proof}

\begin{theorem}[a uniform Bessel majorant on every simplex stratum]
\label{thm:uniform-simplex-majorant}
Let $n_1,\ldots,n_k$ be arbitrary positive integers with sum $N$, and let
$0<u<1$, $v=u/(1-u)$. Define
\[
\mathcal Q_n(u)=(1-u)^{N-1}v^{k-1}
\int_0^\infty e^{-t}t^k\prod_{i=1}^k\Phi(n_i vt)\,dt.
\]
Then, with no lower bound on any $n_i/N$,
\begin{equation}\label{eq:uniform-simplex-majorant}
0\le1-\frac{S_n(u)}{\mathcal Q_n(u)}
\le k(k+1)v+\frac{k}{2}v^2\left(\sum_i\sqrt{n_i}\right)^2
\le k(k+1)v+\frac{k^2}{2}Nv^2.
\end{equation}
Thus this approximation has relative error $O_k((\log N)^2/N)$ whenever
$u=O(\log N/N)$, uniformly over every positive composition of $N$.
\end{theorem}

\begin{proof}
The positive refresh formula gives
\[
S_n(u)=\frac{(1-u)^{N-1}}v
\int_0^\infty e^{-t}\prod_iB_{n_i}(vt)\,dt.
\]
Applying Lemma~\ref{lem:uniform-laguerre-bessel} to each factor, and then
the product-deficit inequality, bounds the relative deficit by
$(kv/2)\mathbb E[T]$, where $T$ has density proportional to
$e^{-t}t^k\prod_i\Phi(n_i vt)$.
For the probability weights proportional to $x^j/[j!(j+1)!]$, the identity
$\mathbb E[J(J+1)]=x$ implies
\[
x\frac{\Phi'(x)}{\Phi(x)}=\mathbb E[J]\le\sqrt{x}.
\]
Integration by parts therefore gives
\[
\mathbb E[T]=k+1+
\sum_i\mathbb E\!\left[\frac{n_ivT\Phi'(n_ivT)}{\Phi(n_ivT)}\right]
\le k+1+\sqrt v\left(\sum_i\sqrt{n_i}\right)\sqrt{\mathbb E[T]}.
\]
The boundary terms vanish, since near zero the density is $O(t^k)$, while
$\Phi(x)\le e^{2\sqrt{x}}$ bounds the tail by a polynomial times
$\exp(-t+2\sqrt{vt}\sum_i\sqrt{n_i})$. This also proves finiteness of
the integral and of every moment used above. Applying $2ab\le a^2+b^2$
to the final term yields
$\mathbb E[T]\le2(k+1)+v(\sum_i\sqrt{n_i})^2$.
The claimed estimates follow, the last one by Cauchy--Schwarz.
\end{proof}

\subsection{Coordinates of order \texorpdfstring{$N/(\log N)^2$}{N divided by (log N) squared}}

A small coordinate $a$ adds its own runs. Each run needs about two switches,
so it gives a factor $u^2$, and it can go in about $N$ places. So what
matters is $aNu^2$, just like $bu^2$ in the binary edge law. At the crossing
$u$ is about $(\log N)/N$, and then $aNu^2$ is about $a(\log N)^2/N$, which
is why coordinates of this size start to matter.

Fix integers $r\ge2$ and $\ell\ge1$, and put $k=r+\ell$.
Consider $r$ large counts $n_1,\ldots,n_r$ and $\ell$ further positive
counts $a_1,\ldots,a_\ell$. Write
\[
N=\sum_i n_i+\sum_j a_j,\quad M=\sum_i n_i,\quad
L=\log N,\quad \beta_i=\frac{n_i}{M},\quad
\rho_j=\frac{a_jL^2}{N}.
\]
The $\beta_i$ and $\rho_j$ below always refer to these actual counts.
Define
\[
A=\sum_{i=1}^r\sqrt{\beta_i},\quad g=A^2-1,\quad
K=r+\ell-1,\quad H=\frac{r-1}{2}+2\ell,\quad c=\frac Kg,
\]
and
\[
D_r(\beta)=\frac{A^{r/2+1}(\prod_i\beta_i)^{-3/4}}
{(4\pi)^{(r-1)/2}}.
\]

\begin{theorem}[a crossover between neighboring simplex strata]
\label{thm:simplex-boundary-crossover}
Fix $\varepsilon>0$ and $R<\infty$. Assume $\beta_i\ge\varepsilon$
and $0<\rho_j\le R$. Uniformly for such count vectors and
$c_0L\le z\le C_0L$, where $0<c_0<C_0<\infty$ are fixed,
\begin{align}
S_{(n,a)}(z/N)
={}&D_r(\beta)A^{2\ell}\frac{z^H e^{gz}}{N^K}
\prod_{j=1}^\ell\Phi\!\left(\frac{A^2a_jz^2}{N}\right)
\left(1+O\!\left(L^{-1}+\frac{L^2}{N}\right)\right).
\label{eq:simplex-boundary-profile}
\end{align}
The unique positive solution $u_{n,a}$ of $S_{(n,a)}(u)=1$ consequently
satisfies
\begin{align}
Nu_{n,a}
={}&cL-\frac Hg\log L
-\frac1g\left\{
\log\bigl(D_r(\beta)A^{2\ell}\bigr)+H\log c
+\sum_{j=1}^\ell\log\Phi(A^2c^2\rho_j)\right\}
\notag\\
&\hspace{20mm}+O\!\left(\frac{\log L}{L}+\frac{L^2}{N}\right).
\label{eq:simplex-boundary-root}
\end{align}
The constants in the errors depend only on the displayed fixed parameters.
\end{theorem}

\begin{proof}
Let $v=u/(1-u)$ and $Y=Mv$. By Theorem~\ref{thm:uniform-simplex-majorant},
with relative error $O_k(v+Nv^2)$, the exact ratio equals
\begin{equation}\label{eq:simplex-mixed-integral}
\frac{(1-u)^{N-1}v^{\ell-1}}{\prod_{i=1}^r n_i}
\int_0^\infty e^{-t}t^\ell
\prod_{i=1}^r\left[\sqrt{\beta_iYt}\,
I_1(2\sqrt{\beta_iYt})\right]
\prod_{j=1}^\ell\Phi(a_jvt)\,dt.
\end{equation}
In the stated regime $Y\asymp L$ and all $a_jvY$ are bounded.
Set $w=\sqrt t$. The exponential from the $r$ large-coordinate factors
is $\exp(A^2Y-(w-A\sqrt Y)^2)$.
On $|w-A\sqrt Y|\le\sqrt Y/2$, the large-coordinate Bessel arguments
are bounded below by a positive constant times $Y$. Their usual
large-argument expansions are uniform, with relative error $O(1/Y)$.
The remaining smooth factor, apart from constants and powers of $Y$, is
\[
w^{r/2+1+2\ell}\prod_j\Phi(a_jvw^2).
\]
Divide this factor by its value at $w=A\sqrt Y$. On the same interval
its second derivative is $O(1/Y)$, uniformly, because $w/\sqrt Y$ lies in a fixed
compact subset of $(0,\infty)$, the quantities $a_jvY$ lie in a fixed
compact interval, and $\Phi\ge1$ there. Taylor's theorem, followed by
integration against the centered Gaussian, eliminates the linear term
and gives relative error $O(1/Y)$.

The part outside this interval is exponentially small, uniformly. The
bounds $I_1(x)\le e^x$ and
$\Phi(x)\le e^{2\sqrt x}$ give a polynomial times
\[
\exp\!\left(A^2Y-(w-A\sqrt Y)^2
+2w\sum_j\sqrt{a_jv}\right).
\]
Here $\sum_j\sqrt{a_jv}=O(Y^{-1/2})$, so completing the square leaves
a Gaussian centered only $O(Y^{-1/2})$ away from $A\sqrt Y$.
On the excluded region its integral is exponentially small relative to
the main term, even after the polynomial factors are included.
The same argument also justifies extending the central Gaussian
integrals to the full real line.

Evaluation of the leading Gaussian integral in
\eqref{eq:simplex-mixed-integral} now gives
\begin{align*}
S_{(n,a)}(u)
={}&D_r(\beta)A^{2\ell}\frac{Y^H}{M^K}
e^{A^2Y+(N-1)\log(1-u)}
\prod_j\Phi(A^2a_jMv^2)\\
&\qquad\times\left(1+O(Y^{-1}+v+Nv^2)\right).
\end{align*}
Since $\sum_j a_j/N=O(L^{-2})$, at $u=z/N$ we have
\[
Y=z+O(L^{-1}+L^2/N),\qquad
A^2Y+(N-1)\log(1-u)=gz+O(L^{-1}+L^2/N).
\]
Replacing $Y$ by $z$, $M$ by $N$, and $A^2a_jMv^2$ by
$A^2a_jz^2/N$ in the other factors changes the expression by a relative
$O(L^{-2}+L/N)$, because the arguments of $\Phi$ stay bounded and
$\Phi\ge1$. This proves~\eqref{eq:simplex-boundary-profile}.

The proper-run polynomial has nonnegative coefficients, zero constant
coefficient, and positive coefficient $k!$ at $u^{k-1}$, so the crossing
is unique. On the compact set of permitted $\beta$, both $g$ and
$D_r(\beta)$ are bounded above and bounded away from zero. Testing
\eqref{eq:simplex-boundary-profile} at fixed sufficiently small and large
multiples of $L$ brackets the root uniformly. Its logarithm gives
\[
gz+H\log z=KL-\log(D_r(\beta)A^{2\ell})
-\sum_j\log\Phi(A^2\rho_j(z/L)^2)
+O(L^{-1}+L^2/N).
\]
First $z=cL+O(\log L)$, so the sum of logarithms on the right differs
from its value at $z=cL$ by $O(\log L/L)$. Substituting this estimate
into $\log z$ proves~\eqref{eq:simplex-boundary-root}. Inverting keeps
the error because the derivative of $gz+H\log z$ is at least $g>0$.
\end{proof}

\begin{corollary}[invisibility of subcritical rare counts]
\label{cor:simplex-rare-invisibility}
Under the hypotheses of Theorem~\ref{thm:simplex-boundary-crossover}, if
all $a_j=o(N/(\log N)^2)$, the sum involving $\Phi$ in
\eqref{eq:simplex-boundary-root} is $o(1)$. Thus, at fixed actual large
composition $\beta$, the rare counts do not affect the expansion through
its constant term. Their number $\ell$ still changes both its leading
and its logarithmic coefficients.
\end{corollary}

The scale $N/(\log N)^2$ applies near a face with at least two coordinates
of order $N$. It does not apply at a binary endpoint, where only one
coordinate is large and $g=0$. That case is treated below and in
Section~\ref{sec:vertex-completion}.
Theorem~\ref{thm:simplex-boundary-crossover} gives no error bound when some
$\rho_j$ is unbounded.

\begin{lemma}[exact insertion of a singleton color]
\label{lem:singleton-insertion}
Let $n$ have total size $M$ and let $\mathcal D=u\,d/du$. Then
\[
S_{(n,1)}(u)=\bigl[2u+(M-1)u^2\bigr]S_n(u)
+(u-u^2)\mathcal D S_n(u).
\]
\end{lemma}

\begin{proof}
Insert the new color into a word of length $M$ with $j$ changes.
The two end positions contribute a factor $2u$, the $j$ change gaps
contribute $ju$, and the $M-1-j$ repeat gaps contribute
$(M-1-j)u^2$. Removing the unique new letter reverses this construction.
Summing the resulting identity over all words proves the formula.
\end{proof}

The singleton identity gives an exact finite-size check of the factor
$A^2z^2/N$ that one very small coordinate contributes in
\eqref{eq:simplex-boundary-profile}.

\subsection{Fixed rare counts near a vertex}

With only one large coordinate, the limit is a finite polynomial instead.
In this result the rare counts are fixed, and the large coordinate $b$,
not the total count, is the asymptotic parameter.

\begin{theorem}[a multicolor Laguerre limit at a vertex]
\label{thm:simplex-fixed-vertex}
Fix $\ell\ge1$ and positive integers $a_1,\ldots,a_\ell$. Define
\[
H(y)=\prod_{i=1}^\ell B_{a_i}(y),\qquad
Q(y)=\sum_{i<j}B'_{a_i}(y)B'_{a_j}(y)
\prod_{h\ne i,j}B_{a_h}(y).
\]
The sum defining $Q$ is zero when $\ell=1$. Uniformly on compact
subsets of $y>0$, as $b\to\infty$,
\begin{equation}\label{eq:simplex-vertex-ratio}
S_{(b,a)}\!\left(\sqrt{y/b}\right)
=H(y)+\frac{2\sqrt y}{\sqrt b}\bigl(H'(y)+yQ(y)\bigr)
+O_a(b^{-1}).
\end{equation}
Let $y_a>0$ be the unique root of $H(y_a)=1$, and put $c_a=\sqrt{y_a}$.
The unique endpoint-crossing odds satisfy
\begin{equation}\label{eq:simplex-vertex-root}
u_{b,a}=\frac{c_a}{\sqrt b}
-\frac{1+y_aQ(y_a)/H'(y_a)}{b}
+O_a(b^{-3/2}).
\end{equation}
For one rare color this recovers the binary fixed-edge correction $-1/b$.
For at least two rare colors the coefficient of $1/b$ is strictly less
than $-1$.
\end{theorem}

\begin{proof}
Fix rare-run counts $m_i\in\{1,\ldots,a_i\}$, put $m=\sum_i m_i$,
and let $q$ be the number of runs of the large color. The rare run lengths
can be chosen in $\prod_i\binom{a_i-1}{m_i-1}$ ways. Since successive
large-color runs must be separated, $q\le m+1$, and the large run lengths
have $\binom{b-1}{q-1}$ choices. At $u=\sqrt{y/b}$, configurations with
these run counts contribute $O_a(b^{(q-m-1)/2})$. There are finitely many
possible $m_i,q$, independently of $b$, so only $q=m+1$ and $q=m$ affect
the expansion through order $b^{-1/2}$.

If $q=m+1$, the run word starts and ends in the large color and alternates
between large and rare colors. There are $m!/\prod_i m_i!$ rare orders.
Using $\binom{b-1}{m}=b^m/m!+O_a(b^{m-1})$ and summing gives $H(y)$,
with an $O_a(b^{-1})$ error.

If $q=m$, there are two possibilities. One endpoint is rare and the run
word otherwise alternates, giving $2m!/\prod_i m_i!$ orders. These
terms sum to $2\sqrt y\,H'(y)/\sqrt b$. Otherwise both endpoints are
large and exactly one pair of consecutive rare runs is present. The
members of that pair must have different colors. For an ordered pair
of colors $i,j$, contracting the marked adjacent pair shows that there
are $(m-1)!/[(m_i-1)!(m_j-1)!\prod_{h\ne i,j}m_h!]$ orders.
Summing over the two orientations and all $i<j$ gives
$2y^{3/2}Q(y)/\sqrt b$. Terms with $q\le m-1$ and the next terms of
the binomial expansions are $O_a(b^{-1})$, uniformly for $y$ in a
positive compact set. This proves~\eqref{eq:simplex-vertex-ratio}.

The polynomial $H$ has positive coefficients, tends to zero at zero,
and grows to infinity, so $y_a$ exists uniquely and $H'(y_a)>0$.
The same positivity argument gives the unique crossing for each $b$.
First bracket it using~\eqref{eq:simplex-vertex-ratio} to obtain
$bu_{b,a}^2\to y_a$. The finite run expansion may also be differentiated
uniformly on this compact neighborhood, and its derivative tends to
$H'(y_a)>0$. Substitution of
$u=c_a/\sqrt b+\alpha/b$ into~\eqref{eq:simplex-vertex-ratio}
then forces $\alpha=-1-y_aQ(y_a)/H'(y_a)$, and the uniform remainder
gives~\eqref{eq:simplex-vertex-root}. If $\ell\ge2$, all terms in $Q(y)$
are positive for $y>0$, proving the strict inequality.
\end{proof}

If all $\ell$ rare colors occur once, then $H(y)=y^\ell$ and
$y_a=1$, giving the explicit specialization
\[
u_{b,1,\ldots,1}=b^{-1/2}-\frac{\ell+1}{2b}
+O_\ell(b^{-3/2}).
\]
An independent exact formula in this case is
\[
S_{(b,1,\ldots,1)}(u)=\ell!
\sum_{q=1}^{\min(b,\ell+1)}
\binom{\ell+1}{q}\binom{b-1}{q-1}u^{q+\ell-1}.
\]
Indeed, order the $\ell$ distinct rare letters, choose the $q$ gaps among
their $\ell+1$ end and interior gaps that contain large-color runs,
and choose the positive lengths of those runs.
In particular two singleton rare colors give
\[
u_{b,1,1}=b^{-1/2}-\tfrac32b^{-1}+O(b^{-3/2}).
\]
This agrees with the exact polynomial
$S_{(b,1,1)}(u)=6u^2+6(b-1)u^3+(b-1)(b-2)u^4$.
Theorem~\ref{thm:simplex-fixed-vertex} is not uniform in the rare counts.
Section~\ref{sec:vertex-completion} gives the leading crossing uniformly
for total rare count $A=o(b)$.
